\section{Introduction}
\label{sec:introduction}
Representation learning is usually assessed through the performance of fitted predictors. If a predictor misses a target accuracy, the observation alone does not show whether the representation is inadequate, the allowed decoder class is too restrictive, or the training procedure failed to access an adequate function. This distinction matters whenever a model must meet an error tolerance under a declared resource or sensitivity constraint.

We study a fixed observed-output prediction task. A representation $\Phi$ maps an available history $X$ to coordinates $U$, and a declared metric and sensitivity budget specify the admissible decoders. A lower bound on the best risk in that class can rule out every admissible decoder. An upper bound on the risk of one admissible fitted predictor can establish adequacy. When every cheaper declared class is excluded and a fitted member of another class is accepted, the latter has minimum feasible cost within that family. A positive lower bound alone does not identify removable information loss: irreducible target noise and sensitivity restrictions also contribute.

The main construction strengthens a pairwise obstruction using observation blocks. Within each block we minimize squared error subject to all pairwise Lipschitz constraints, and normalize by one less than the block size. For scalar targets, the resulting expected scores increase with block size and approach optimal Lipschitz risk under the stated boundedness, compactness, and unconditional iid assumptions. This population result does not imply that larger blocks yield better finite-sample endpoints: each larger block consumes observations and reduces the number of independent scores available for calibration.

The paper separates three forms of assurance (Table~\ref{tab:assurance}). They share mathematical and computational tools but answer different questions. In particular, a numerical packet confirming a reported output is not a certificate that an unknown future response lies near that output.

\begin{table}[t]
\centering\small
\begin{tabular}{@{}p{.21\linewidth}p{.36\linewidth}p{.37\linewidth}@{}}
\toprule
Object & Supported conclusion & Essential premise\\
\midrule
Population class risk & Every decoder in a declared class exceeds a tolerance, or a fitted member meets it & Fixed class and loss; justified sampling; valid simultaneous calibration\\
Finite-cohort class risk & Exclusion or constructive adequacy on exactly the specified observations & Fixed coordinates, targets and class; sound optimization witnesses\\
Numerical reported output & Reported value lies within a numerical tolerance of a specified reference computation & Trusted feature input, fixed head semantics, valid arithmetic and identity checks\\
\bottomrule
\end{tabular}
\caption{The three uses of a certificate remain separate throughout the paper. No row implies the premises or conclusion of another.}
\label{tab:assurance}
\end{table}

Our contributions are organized around these questions.
\begin{enumerate}[leftmargin=*,itemsep=2pt]
\item \textbf{A centered scalar risk hierarchy.} We connect the exact pair score to larger constrained block minima, prove monotonicity and consistency in the stated setting, and characterize finite-block limitations. A resolution bound separates approximation bias, statistical uncertainty, and numerical slack.
\item \textbf{Computable and interpretable decisions.} Verified lower witnesses, independent calibration, fitted-risk upper bounds, and retained-block geometry support class exclusion, adequacy, and family-relative minimum-cost statements. Recycled block duals supply a second calibration without additional optimization solves. Standard concentration, scalar extension, and weak duality remain attributed foundations.
\item \textbf{Evidence about when the conclusions help.} Controlled decisions and a natural finite-cohort exclusion are complemented by a class-versus-head diagnosis, calibration limits, and unsuccessful screening and deployment studies. A separate numerical output check is evaluated on frozen forecasting policies with strong conventional controls.
\end{enumerate}

KAHM dictionaries are an implementation case study rather than a premise of the risk theorem. Their reconstruction affinities can be retained and expanded while preserving earlier feasible decoders. Raw-history and Nystr\"om representations supply essential controls. Scaling, local subspace cap, representation metric, and decoder budget are explicit parts of the candidate specification. The accumulated evidence does not establish general KAHM superiority or certificate-directed savings from avoiding model fits.

